\begin{figure}[t]
\centering
% Shared color TikZ styles for the paper figures.
\colorlet{paperBlue}{blue!9!white}
\colorlet{paperBlueLine}{blue!55!black}
\colorlet{paperTeal}{teal!10!white}
\colorlet{paperTealLine}{teal!55!black}
\colorlet{paperAmber}{orange!15!white}
\colorlet{paperAmberLine}{orange!60!black}
\colorlet{paperRose}{red!9!white}
\colorlet{paperRoseLine}{red!55!black}
\colorlet{paperViolet}{violet!10!white}
\colorlet{paperVioletLine}{violet!55!black}
\colorlet{paperInk}{black!72}

\tikzset{
  paperfig/.style={
    font=\small,
    >=Stealth,
    line cap=round,
    line join=round
  },
  paperline/.style={
    -{Stealth[length=2.1mm,width=1.6mm]},
    line width=0.58pt,
    draw=paperInk
  },
  paperdash/.style={
    paperline,
    dashed,
    draw=black!48
  },
  paperbox/.style={
    draw=black!58,
    line width=0.58pt,
    rounded corners=2pt,
    fill=white,
    align=center,
    inner xsep=5pt,
    inner ysep=5pt,
    minimum height=8mm
  },
  paperdata/.style={
    paperbox,
    fill=paperBlue,
    draw=paperBlueLine
  },
  paperproc/.style={
    paperbox,
    fill=paperTeal,
    draw=paperTealLine
  },
  papermodel/.style={
    paperbox,
    fill=paperAmber,
    draw=paperAmberLine,
    line width=0.65pt
  },
  paperout/.style={
    paperbox,
    fill=paperRose,
    draw=paperRoseLine,
    line width=0.65pt,
    font=\small\bfseries
  },
  paperstate/.style={
    paperbox,
    fill=paperViolet,
    draw=paperVioletLine
  },
  papernote/.style={
    font=\scriptsize,
    align=center,
    text=black!70,
    fill=white,
    inner sep=1.4pt
  },
  paperaxis/.style={
    line width=0.45pt,
    draw=black!72
  },
  papergrid/.style={
    line width=0.25pt,
    draw=black!15
  },
  paperplot/.style={
    line width=0.95pt,
    draw=paperBlueLine
  },
  paperplotA/.style={
    line width=0.95pt,
    draw=paperBlueLine
  },
  paperplotB/.style={
    line width=0.95pt,
    draw=paperAmberLine
  },
  papermark/.style={
    circle,
    draw=paperBlueLine,
    fill=white,
    line width=0.55pt,
    inner sep=1.6pt
  },
  papermarkA/.style={
    papermark,
    draw=paperBlueLine,
    fill=paperBlue
  },
  papermarkB/.style={
    papermark,
    draw=paperAmberLine,
    fill=paperAmber
  }
}

\resizebox{0.98\linewidth}{!}{%
\begin{tikzpicture}[paperfig, node distance=7mm and 8mm]
  \node[paperdata, text width=20mm] (task) {当前任务\\$x_i$};
  \node[papermodel, right=of task, text width=22mm] (gen) {候选生成\\$y_i^{(1)}$};
  \node[paperproc, right=of gen, text width=22mm] (exec) {统一执行器\\测试判定};
  \node[paperproc, below=of exec, text width=22mm] (classify) {失败分类\\错误标签};
  \node[paperproc, left=of classify, text width=22mm] (compress) {反馈压缩\\$\phi(e_i^{(1)})$};
  \node[paperproc, left=of compress, text width=20mm] (repair) {修补/替换\\$\tilde{x}_i^{(2)}$};
  \node[paperout, below=of classify, text width=22mm] (stop) {通过或停止};

  \draw[paperline] (task) -- (gen);
  \draw[paperline] (gen) -- (exec);
  \draw[paperline] (exec) -- node[papernote, right] {失败} (classify);
  \draw[paperline] (exec.east) to[bend left=24]
    node[papernote, right] {通过} (stop.east);
  \draw[paperline] (classify) -- node[papernote, right] {预算耗尽} (stop);
  \draw[paperline] (compress) -- (repair);
  \draw[paperline] (classify) -- (compress);
  \draw[paperline] (repair.north) to[bend left=22]
    node[papernote, left] {下一轮} (gen.south);

  \node[papernote, below=6mm of repair, text width=54mm] (note) {仅回注局部约束，控制上下文膨胀};
  \draw[paperdash] (compress.south) -- (note.north);
\end{tikzpicture}%
}
\caption{失败反馈救援闭环}
\label{fig:rescue_loop}
\end{figure}
